\newpage
\section{Conclusion and Next Steps}

\subsection{Semester 2 Aims}
With an understanding of all IMU operations, a baseline of how the UKF is able to fuse
sensor measurements and its performance, and with the residual block of the TCN implemented,
semester two is on track for the completion of this project. The main core of the model is the 
residual block and from here the rest of the work is covered in the next subsections.

\subsubsection{Model Hyperparameters}
The main thing to do is to determine multiple of the model hyperparameters such as, the filter
size, recpetive field, dilation factor, and the number of channels per residual block. For 
dilation factor, it is a common method to follow is $d = 2^i$ or $d = 4^i$ where $i$, 
is the layer number as seen by \cite{brossard_2020_denoising}. This is because by increasing
the dilation factor the model can see further back in time and use this information to more
accurately predict the accumulation of errors. However, as this is not a standard rule a wide
variety of dilation factors will be used to test model performance. 




\subsubsection{Model Training}


\subsection{Summary of Progress}
During Semester One, the project progressed from initial understanding 
of the challenges faced in orientation estimation, to exploration and 
designing of deep learning models. Much of the work conducted
had been research using academic literature and online
lectures. With foundational knowledge gained in all areas of the project,
it has ensured that Semester Two can be focused on the completion    
deep learning model to address the questions posed in
\hyperref[sec:question]{section 1.3}. Furthermore, the exploration of how the 
baseline UKF performed also gave me vital understanding in how the model should attempt
to learn gyroscopic drift.

\subsection{Challenges}
The primary challenge encountered during this semester stemmed from having
no prior knowledge of IMU operations and deep learning. This led to 
extensive research needed for the completion of the project, 
which slowed down the implementation of the model. However, now having
concluded research I have been able to make informed decisions behind 
the selection of model, optimisation through normalisatio, and measurements
that accurately assess the performance of the model. 

\subsection{Skills Developed}
The main skill developed has been learning how to assess and evaluate 
literature and ascertain the correct conclusions to be able to be applied
to this project. My ability to use the PyTorch library and 
Python skills have developed to a significant degree in which modelling,
measurements, and analysis will be conducted within Python. In addition, 
my technical writing capabilities have also progressed especially compared
to my third-year project. 
